\chapter*{Resumen}
\addcontentsline{toc}{chapter}{Resumen}
\label{sec:resumen}

La caza de amenazas sobre grafos heterogéneos construidos a partir
de logs de seguridad choca contra tres obstáculos recurrentes:
(i) ingesta frágil frente a formatos nuevos o mutantes, (ii) costo
combinatorio del \emph{matching} de patrones en grafos grandes con
múltiples tipos de entidades y aristas, (iii) señales de amenaza
mezcladas con ruido operacional sin un esquema de priorización
estable. Los SIEM tradicionales responden sobre tablas con
lenguajes orientados a texto; el analista pierde la estructura
relacional del incidente en el momento mismo de formular la
pregunta.

Este reporte describe GraphHunter, una arquitectura diseñada para
atacar los tres obstáculos por separado. La contribución principal
es una descomposición en \emph{shearing layers} que aísla tres
capas con velocidades de cambio distintas: un núcleo estable con el
motor de grafo, el matcher SIMD y el gate de ML; una plataforma
con superficies (HTTP, MCP, DSL, parsers, canonical OCSF); y un
conjunto de aplicaciones delgadas. Sobre esa descomposición
presentamos: un matcher de patrones con intersección AVX2 de
posting-lists cuya equivalencia con una referencia DFS se verifica
por property-testing; un DFS temporal guiado por score para la fase
de \emph{refinement}; un compositor de anomalía de cinco
componentes con pesos operables; un aislamiento FFI para el motor
C++ mediante \emph{typestate} y RAII con versionado ABI explícito;
un \emph{model gate} que hace opcional al GNN de amenazas sin
degradar el camino base; y un ingestor agéntico de tres carriles
(VRL compilado, LLM-como-compilador offline, drift-como-señal) que
combina determinismo productivo con adaptación offline gobernada.
La evaluación incluye benchmarks de latencia y throughput,
property-tests de equivalencia DFS↔SIMD, y un smoke manual
end-to-end sobre datasets sintéticos. Concluimos con limitaciones
abiertas: dos aristas residuales \code{core}$\to$\code{platform} y
la ausencia de un modelo GNN entrenado sobre FortiAnalyzer
producción.

El libro complementario \emph{Diseño completo}
(\filepath{docs/diseno/diseno.pdf}) profundiza en el diseño
algorítmico del \emph{matcher} y los índices auxiliares con
diagramas pedagógicos: el \emph{stack} de operadores físicos
(DFS, LFTJ-K4, Yannakakis chain, Yannakakis tree-walk), la
\emph{lazy build} y el \emph{sparse mmap-spill} de NLF y k-hop
reachability, y el plano de obra de las stages futuras de la
\emph{streaming layer}.
